\item Las paredes de la celda de una prisión están orientadas en las cuatro direcciones cardinales. En el primer día de primavera, la luz del Sol naciente entra a una ventana rectangular en la pared oriente. La luz recorre $2.37\text{ m}$ horizontalmente para aparecer en forma perpendicular en la pared poniente opuesta.
\begin{enumerate}
    \item ¿Con qué rapidez se mueve el rectángulo iluminado en la pared poniente?
    \item El prisionero sostiene un pequeño espejo cuadrado contra la pared en una esquina del rectángulo de luz, reflejándola hacia su origen en la pared oriente. ¿Con qué rapidez se mueve el cuadrado de luz de menor tamaño en esa pared?
    \item Visto desde una latitud de $40.0^\circ$ al norte, el Sol naciente se traslada en el cielo a lo largo de una línea que forma un ángulo de $50.0^\circ$ con el horizonte del sureste. ¿En qué dirección se mueve el rectángulo de luz en la pared poniente?
    \item ¿En qué dirección se mueve el cuadrado de luz más pequeño en la pared oriente?
\end{enumerate}